% =====================================================================
%  GLOSARIO DE TÉRMINOS
% =====================================================================
\chapter*{Glosario de Términos}
\addcontentsline{toc}{chapter}{Glosario de Términos}

\begin{singlespace}
\renewcommand{\arraystretch}{1.3}
\begin{longtable}{@{}L{4cm} L{11.4cm}@{}}
Acoso sexual & Conducta de naturaleza sexual no deseada que afecta la dignidad de la persona. En la normativa peruana se denomina hostigamiento sexual. \\
Adjudicación & Revisión de un especialista que resuelve los desacuerdos entre anotadores y fija la etiqueta final de un relato. \\
Ajuste fino & Entrenamiento adicional de un modelo preentrenado con datos etiquetados de una tarea específica (\textit{fine-tuning}). \\
Anonimización & Proceso de retirar o reemplazar la información que permite identificar a una persona. \\
Anotación & Asignación de una etiqueta a cada relato siguiendo una guía con criterios explícitos. \\
Corpus & Conjunto organizado de textos que se usa para analizar y entrenar modelos. \\
Desbalance de clases & Situación en la que una clase tiene muchos más ejemplos que otra. \\
Exhaustividad & Proporción de los casos positivos reales que el modelo logra identificar (\textit{recall}). \\
Fuga de información & Uso, durante el entrenamiento o la selección del modelo, de datos que debían reservarse para la evaluación. \\
Línea base & Modelo sencillo que sirve de punto de comparación para el modelo desarrollado. \\
Línea base léxica & Línea base que marca como acoso todo relato que contiene alguna palabra de una lista de términos sexuales. \\
Precisión & Proporción de los casos que el modelo marcó como positivos que realmente lo son. \\
Prototipo & Versión funcional de una aplicación que permite probar su utilidad antes de un desarrollo definitivo. \\
Relato & Texto que cuenta una experiencia, con referencias a lo que ocurrió, su contexto o sus consecuencias. \\
Sobreajuste & Situación en la que un modelo memoriza los datos de entrenamiento y falla con datos nuevos. \\
Validación cruzada & Procedimiento que divide los datos de desarrollo en pliegues para entrenar y validar el modelo varias veces. \\
\end{longtable}
\end{singlespace}
\setcounter{table}{0}
